\chapter{Conclusiones y Líneas Futuras}
\label{cap:conclusiones}

\section{Resumen}

Esta memoria ha descrito un método para predecir la llegada de autobuses a sus paradas en una red cuyo operador
publica la posición de sus vehículos pero ningún retraso, y cuyo horario intermedio es una interpolación a velocidad
constante. El método se resume en cuatro pasos:

\begin{enumerate}
  \item \textbf{Situar} el GPS sobre el trazado oficial, en una ventana alrededor de la próxima parada que declara el
        feed y con una regla de desempate para los recorridos que vuelven sobre sí mismos, y traducir esa posición a
        la hora a la que el horario pone allí al autobús.
  \item \textbf{Sumar} el tiempo real que cuesta recorrer lo que falta, tramo a tramo, con tiempos aprendidos por
        franja, contraídos hacia el tiempo del tramo a cualquier hora, y con una cascada de respaldo por longitud y por
        factor de ritmo.
  \item \textbf{Elegir} en cada predicción la mejor fuente fiable entre el GPS, la última parada registrada y el
        horario corregido, con reglas explícitas para abstenerse y sin afirmar nunca un retraso que no se pueda
        sostener.
  \item \textbf{Aprender y medir} cada día con el propio feed, y cambiar los parámetros solo si el error baja en
        recorridos que no se usaron para aprenderlos; las reglas nuevas, solo si mejoran en días no vistos al aprender.
\end{enumerate}

\section{Resultados}

En 2,21~millones de llegadas a menos de veinte minutos, recogidas en 17 días, la llegada mostrada cae dentro de $\pm2$~minutos
en el 96\,\% de los casos cuando el autobús está a menos de tres minutos y en el 67\,\% entre diez y veinte minutos,
frente al 24\,\% y el 25\,\% del horario publicado, con intervalos de confianza de, como mucho, tres puntos. En una muestra
de 14~minutos con GPS, la llegada no tiene sesgo y acierta el 94\,\% a menos de tres minutos. Sin GPS, con solo la
parada registrada, el acierto es el mismo a menos de tres minutos (94\,\%) y baja seis puntos entre diez y veinte
minutos (del 73 al 67\,\% en esa muestra). Sin ningún dato del
vehículo, el horario corregido duplica o triplica el acierto del horario publicado, y la corrección por viaje baja su
error medio de 210 a 158~s; con el percentil 40 de los desvíos en lugar de la media, la hora sin autobús acierta
entre 2 y 4 puntos más a $\pm2$~minutos en días no vistos al aprender. Cuando todo el camino tiene tiempos
aprendidos, el 93\,\% de las llegadas a menos de diez minutos cae dentro de $\pm2$~minutos, el techo que permite
para ellas la resolución de este feed. Con el protocolo de evaluación por días, un filtro de Kalman sobre el ritmo del propio autobús y
una red neuronal reducen el error de las sumas de tramos un 4\,\% y entre un 6 y un 7\,\%, respectivamente.

\section{Lecciones}

\subsection{Lo Único Vivo Es la Posición} Antes de modelar nada hubo que comprobar qué datos eran fiables. El retraso
declarado no existía y el horario intermedio era una interpolación. Todo el diseño parte de usar la posición y
diferencias de horario, que hacen irrelevante el retraso acumulado (proposición~\ref{prop:diferencias}).

\subsection{Revisar los Datos de Entrada} Durante semanas, 670 de los 6\,130 viajes (el 11\,\%, en 25 líneas) no se
pudieron situar con el GPS porque el horario oficial dibujaba sus trazados al revés, y esos viajes tampoco aprendían.
No era un fallo del método, sino de los datos, y solo se vio al revisar uno a uno los autobuses que no salían en
directo. Corregir los trazados en la construcción y situar al autobús que espera en la cabecera subió los autobuses
situados del 73{,}8 al 97{,}4\,\%, más que cualquier mejora del modelo.

\subsection{Medir Antes de Adoptar} Muchas ideas razonables empeoraron la medida: corregir con la velocidad del
último minuto,
mezclar lo aprendido con el horario, tratar los festivos como fin
de semana en los tiempos por tramo, suponer que el autobús de la cabecera sale a su hora, arrastrar el retraso del
viaje anterior o seguir al autobús que desaparece del feed; calcular la franja de cada tramo a la hora prevista de paso no
cambió nada. En la capa 3 corregida, en cambio, tratar los
festivos como fin de semana sí mejora, mientras que la primera versión de la corrección, que los trataba por su día de
la semana, los dejaba peor que sin corregir: cada caso se midió por separado. La regla de adoptar solo lo que mejora la medida, por días y de punta a punta, ha sido más útil que
cualquier intuición sobre el tráfico.

\subsection{Abstenerse Es Parte de la Precisión} Mostrar «Estimado» o «Solo horario» cuando no hay datos, ocultar un
retraso no creíble y desconfiar de dos fuentes que se contradicen no aumentan los aciertos medidos, pero evitan los
errores que más daño hacen a la confianza del usuario.

\subsection{Lo Sencillo y Auditable Compite} Medianas por tramo, una contracción con una sola constante, una cascada
de cuatro niveles: cada parámetro tiene un significado físico y se puede comprobar a mano. Frente a una red neuronal
entrenada con lo mismo y con lo que pasa cada día, se pierde entre un 6 y un 7\,\% de error medio; a cambio, cualquier
cifra que muestre la aplicación se puede explicar.

\section{Líneas Futuras}

\begin{itemize}
  \item \textbf{El ritmo del propio autobús.} El filtro de Kalman sobre la desviación del propio autobús mejora las
        sumas de tramos un 4\,\% con dos parámetros (sección «\nameref{sec:comparacion}»). El paso siguiente es medirlo
        sobre la llegada completa, con la regla de adopción de la sección «\nameref{sec:eval-contraccion}».
  \item \textbf{Capas a lo largo de los días.} La captura continua guarda desde el 31 de octubre de 2026 los errores de
        cada bloque por capa, horizonte, franja y tipo de día. Con semanas de datos, las Tablas~\ref{tab:capas}
        y~\ref{tab:cobertura} se podrán repetir con intervalos de confianza estrechos.
  \item \textbf{Temporada y tipo de día.} Con meses de muestras se podrá medir si conviene separar el verano y las
        vacaciones escolares, o el sábado del domingo, en la corrección por viaje y en los tiempos por tramo.
  \item \textbf{Fechas especiales.} Los tiempos aprendidos distinguen laborable y fin de semana, pero no los días de
        fiesta, de eventos o de Navidad. Marcar esas fechas para que no cuenten al aprender evitaría que contaminen los días
        normales; aprenderlas aparte exige capturas del año anterior.
  \item \textbf{Viajes sin autobús.} Distinguir en la aplicación el viaje que ya debería haber salido y no tiene
        autobús en el feed (en torno al 4\,\% de los que deberían estar en marcha), en lugar de darlo como estimado.
        Antes de su hora, la ficha ya dice que aún no tiene autobús asignado.
  \item \textbf{De madrugada.} De 23:58 a 05:00 el feed de Bizkaibus no se actualiza y los nocturnos salen sin directo
        (sección «\nameref{sec:mas-de-24h}»). Si el operador lo mantiene activo de noche, el emparejador de ayer los situará sin
        cambios; para medir esas llegadas y aprender de ellas, la captura tendría que separar también los viajes de ayer.
  \item \textbf{La secuencia de cada vehículo.} En una cabecera, el autobús que empieza un viaje suele ser uno que
        acaba de terminar otro, y su retraso pasa al siguiente. La clave del feed no lo dice (solo el 11\,\% de los
        viajes consecutivos de una clave los hace el mismo autobús), pero las capturas guardan qué vehículo hace cada viaje
        cada día. Si esa secuencia se repite de un día a otro, se podría seguir al autobús en su viaje anterior y
        estimar la salida real; hay que medir antes si se repite.
  \item \textbf{Revisión externa.} Que otras personas repliquen el método con la implementación de referencia y lo
        apliquen a otra red.
  \item \textbf{Intervalos de predicción.} Las muestras permiten estimar, además de la mediana de cada tramo, su
        dispersión. Sumando varianzas por tramo se podría mostrar un intervalo («llega en 6--8 minutos») con una
        cobertura medida.
  \item \textbf{Más capacidad técnica.} Un servidor propio y una base de datos con años de historia quitarían los
        topes de la sección «\nameref{sec:infraestructura}»: aprender festivos, fiestas y temporadas con el año
        anterior, cruzar meteorología y eventos, recalibrar cada hora y servir modelos más capaces. Hay que medir
        cuánto mejora cada paso.
  \item \textbf{Más fuentes de datos.} Pasar del 93\,\% en las llegadas a menos de diez minutos con todo el camino
        aprendido exige información que este feed no tiene: posiciones
        más frecuentes, apertura de puertas o conteo de viajeros.
  \item \textbf{Incidencias.} Los cortes y desvíos publicados por el operador no entran hoy en el cálculo hasta que se
        notan en la posición. Incorporarlos exige medir antes cómo afectan.
  \item \textbf{Planificación de viajes.} Las llegadas estimadas pueden alimentar un buscador de rutas con transbordos
        (por ejemplo, con el algoritmo RAPTOR), de modo que el viaje propuesto use las horas reales y no las
        publicadas.
  \item \textbf{Más redes.} Aplicar el método a otras ciudades y comprobar cuánto de lo medido aquí se mantiene.
\end{itemize}

\section{Reproducibilidad}

Para que cualquiera pueda reproducir el método, esta memoria incluye todas las constantes
(Apéndice~\ref{ap:constantes}), el pseudocódigo (Apéndice~\ref{ap:algoritmos}), los formatos de datos
(Apéndice~\ref{ap:formatos}), vectores de prueba generados con el código de producción (Apéndice~\ref{ap:vectores}) y
una implementación de referencia independiente que los reproduce exactamente (Apéndice~\ref{ap:referencia}). Una
implementación nueva, en cualquier lenguaje, puede darse por correcta cuando reproduce todos los vectores al segundo.
